% Sección 10: Conclusiones y trabajo futuro

\section{Conclusiones y trabajo futuro}

MediFlow entrega una base analítica reproducible e interpretable para estudiar la demanda y la presión operativa del Hospital Universidad del Norte a partir de agregados compatibles con OMOP CDM v5.4, articulando comprensión del flujo histórico, predicción con evaluación temporal estricta y exploración paramétrica de escenarios bajo supuestos declarados. La verificación empírica en el corte documental del repositorio (\texttt{803c1a1}) confirma un contraste riguroso entre fuentes: en el experimento puntual LightGBM Poisson con nivel local (\texttt{lgbm\_final}), el WAPE diario global desciende de 18,0\,\% a 10,5\,\% (mejora relativa archivada de 41,6\,\%) con sólido desempeño en Urgencias (8,4\,\%) y Rehabilitación laboral, mientras que el artefacto cuantílico efectivamente servido en la web (\texttt{v2.0}) exhibe WAPE de 8,91\,\%, 11,08\,\% y 42,18\,\% por servicio con coberturas de 78,57\,\%, 64,29\,\% y 73,08\,\% sobre 182 días. En Hospitalización, la verificación documentada de stock inicial nulo en 12 cortes y de Ley de Little sobre población común depurada por $p_{99{,}9}$ (\num{277,52} pacientes frente a \num{277,47} de media horaria) aporta consistencia interna sólida, al tiempo que la simulación de adelanto del 35\,\% de altas vespertinas estima una reducción relativa de 32,8\,\% en horas rojas (9,32\,\% a 6,26\,\%) sin reducir el pico máximo de 359 pacientes.

Estos resultados retrospectivos justifican avanzar hacia una validación operativa supervisada, pero no demuestran beneficio clínico ni sustituyen las mediciones físicas del hospital. Quedan priorizadas cinco líneas de trabajo inmediato antes de un despliegue institucional: (1) regenerar el artefacto web hacia la versión \texttt{v2.1} desde agregados autorizados para habilitar \texttt{perfiles\_horarios} reales por corte y unificar Ley de Little sobre la población alineada; (2) recalibrar los intervalos cuantílicos e incorporar la selección completa de factores externos aceptados y nivel local dentro del paquete servido, reconciliando además las métricas de largo plazo entre el reporte histórico y el artefacto; (3) abrir, revisar visualmente y republicar el proyecto PBIP de Power BI aclarando o unificando la brecha de corte entre abril y julio de 2026; (4) auditar los sellos temporales de transición e instrumentar las tres marcas del HIS junto con capacidad física y tiempos reales de atención; y (5) ejecutar un piloto en modo sombra y completar el protocolo de borrado seguro de datos locales al término del certamen.
